\section{Towards an NPE-based Shear Estimator}
\label{sec:shear}
\XW{Same model and test set as Sec.~\ref{sec:results} (100k training set, Xu3 only).}

In Sec.~\ref{sec:results}, we showed that our NPE produces well-calibrated posteriors. However, what we ultimately need is an estimate of the shear that is unbiased over an ensemble of galaxies, i.e. one with small $m$ and $\mathbf{c}$ in Eq.~(\ref{eqn:linear-bias}). In this section, we explore how the NPE posteriors can be turned into such an estimate. We start from the naive choice, the posterior mean of each galaxy, and show that it is strongly biased (Sec.~\ref{ssec:raw-bias}). We explain where this bias comes from in Sec.~\ref{ssec:shrinkage}, and then explore several ways to reduce it: weighting galaxies and cutting on their posterior widths (Sec.~\ref{ssec:weights-cuts}), calibrating the shear response (Sec.~\ref{ssec:response}), and, most importantly, combining the full likelihoods of galaxies that share the same shear instead of using point estimates (Sec.~\ref{ssec:stacking}). We present these as explorations rather than as a final, calibrated shear estimator framework, which we leave to future work.

\subsection{Raw Posterior Means}
\label{ssec:raw-bias}

For each test galaxy $i$, we take the mean of its TF-weighted posterior, $\hat{g}_{k,i}$, as a point estimate of its shear. We then measure $m$ and $c$ separately for the two shear components with a linear least-squares fit of $\hat{g}_{k,i} - g_{k,i} = m_k\, g_{k,i} + c_k$, which is Eq.~(\ref{eqn:linear-bias}) with $\hat{g}$ in place of $\hat{\epsilon}$. We use all Xu3 galaxies with equal weights and no cuts, and quote the standard errors of the fit as our uncertainties. We restrict the fit to galaxies with $|g_k| < 0.02$, since this range covers the shear of the vast majority of galaxies in cosmic shear measurements. At larger shears, the response of our estimator becomes slightly nonlinear, with a term cubic in the shear. Such third-order shear biases have also been found for other shear estimators, and become important at the large shears near galaxy clusters \citep{Anacal, AnaCalCluster}.

\begin{figure}
    \includegraphics[width=\linewidth]{figs/fig_raw_bias.pdf}
    \caption{Bias of the raw posterior means on the Xu3 test set ($10^5$ galaxies, equal weights, no cuts). Points show the mean of $\hat{g}_k - g_k$ in bins of the true shear for $g_1$ (blue circles) and $g_2$ (red squares); error bars are smaller than the markers. Solid lines are the linear fits over $|g_k| < 0.02$ (shaded), and dashed lines extend them to the full prior range.}
    \label{fig:raw-bias}
\end{figure}

Fig.~\ref{fig:raw-bias} shows the mean of $\hat{g}_k - g_k$ in bins of the true shear. The additive bias is consistent with zero for both components, but the multiplicative bias is large, with $m_1 = -0.616 \pm 0.014$ and $m_2 = -0.592 \pm 0.014$. In other words, the posterior means respond to only about 40\% of the true shear. Multiplicative biases of tens of percent are not unusual for uncalibrated shear estimators. In DES Y1, for example, the uncalibrated bias of \textsc{im3shape} ranges from $m\approx-0.1$ for high-SNR galaxies to $m\approx-0.3$ for the lowest-SNR galaxies \citep{DESY1Shape}, with similar levels found for HSC \citep{HSCY1Sims}. These biases are then calibrated with image simulations, which motivates the calibration we explore in Sec.~\ref{ssec:response}. Finding a significant multiplicative bias in the point estimate also does not contradict the calibration results of Sec.~\ref{ssec:calibration}. As we show in Sec.~\ref{ssec:shrinkage}, a bias of this kind is expected for posterior means when the data constrain the shear only weakly compared to the prior, even if the posteriors are correct.

We also measure the shape noise of each galaxy following \citep{KLDESI}, as $\sigma_{\epsilon,i} = \sqrt{\sigma_{g_1,i}^2 + \sigma_{g_2,i}^2}$, where $\sigma_{g_k,i}$ is the standard deviation of the TF-weighted posterior of $g_k$. Across the Xu3 test set, the median $\sigma_{\epsilon,i}$ is 0.066, with 68\% of galaxies between 0.040 and 0.079. About a third of the galaxies have $\sigma_{\epsilon,i} > 0.075$, close to the value of $\sqrt{2}\sigma_\pi = 0.082$ that we would get if the posteriors of both components were equal to the prior. Here $\sigma_\pi = 0.2/\sqrt{12} = 0.058$ is the standard deviation of the uniform prior $\mathcal{U}(-0.1, 0.1)$ on each shear component, and the factor of $\sqrt{2}$ comes from adding the two components in quadrature. The data of these galaxies carry little information about the shear, and as we will see in the next section, they are the main source of the bias in the posterior means. We compare our shape noise with that of \citep{KLDESI} in Sec.~\ref{sec:discussion}.

\subsection{Shrinkage of Posterior Means}
\label{ssec:shrinkage}

The bias of the posterior means can be understood as follows. Consider a single shear component $g$ of one galaxy, and suppose that its data give a Gaussian likelihood of width $\sigma_L$, centered on a noisy measurement $g_\mathrm{obs}$ whose average over noise realizations is the true shear. If the prior on $g$ is also a Gaussian centered on zero with width $\sigma_\pi$, the posterior is a Gaussian with mean and variance
\begin{equation}
\label{eqn:gauss-posterior}
    \hat{g} = \frac{\sigma_\pi^2}{\sigma_\pi^2 + \sigma_L^2}\, g_\mathrm{obs}, \qquad
    \sigma_\mathrm{post}^2 = \frac{\sigma_L^2\, \sigma_\pi^2}{\sigma_L^2 + \sigma_\pi^2}.
\end{equation}
The posterior mean is pulled toward the center of the prior by an amount that depends on how informative the data is compared to the prior. Averaging over noise realizations at a fixed true shear $g$ gives $\langle \hat{g} \rangle = g\,\sigma_\pi^2/(\sigma_\pi^2 + \sigma_L^2)$, so using Eq.~(\ref{eqn:gauss-posterior}) to write the result in terms of the posterior width, the multiplicative bias is
\begin{equation}
\label{eqn:shrinkage}
    m = -\frac{\sigma_L^2}{\sigma_L^2 + \sigma_\pi^2} = -\frac{\sigma_\mathrm{post}^2}{\sigma_\pi^2}.
\end{equation}
A galaxy with $\sigma_L \ll \sigma_\pi$ has $m \approx 0$, while a galaxy whose data carry no shear information has $\sigma_\mathrm{post} \to \sigma_\pi$ and $m \to -1$. For a sample of galaxies, the fitted $m$ is approximately the average of the values of the individual galaxies, $m \approx -\langle \sigma_\mathrm{post}^2 \rangle / \sigma_\pi^2$. Our shear prior is uniform rather than Gaussian, so we approximate it by a Gaussian with the same variance, $\sigma_\pi^2 = 0.2^2/12$.

This bias is a natural result from the posterior mean, and would be present even for exact posteriors. For correct posteriors, the average true shear of galaxies with a given $\hat{g}$ is equal to $\hat{g}$, but the average $\hat{g}$ of galaxies with a given true shear, which is what $m$ measures, is shrunk toward zero.

To test whether shrinkage explains the bias we measured in Sec.~\ref{ssec:raw-bias}, we compare Eq.~(\ref{eqn:shrinkage}) with the $m$ measured in groups of galaxies with different posterior widths. We split the test set into 15 equal-size bins of $\sigma_{\epsilon,i}$, and in each bin fit $m_k$ as in Sec.~\ref{ssec:raw-bias} and compute the average posterior variance $\langle \sigma_{\mathrm{post},k}^2 \rangle$. Note that grouping galaxies by their measured posterior widths is not the most stringent version of this test. The posterior width of a galaxy depends on its measured shear, with galaxies that appear more strongly sheared tending to have narrower posteriors. Selecting galaxies with narrow posteriors therefore preferentially selects galaxies whose measured shear $|\hat{g}|$ is larger than their true shear. This selection effect is strongest in the bins with the narrowest posteriors.

\begin{figure}
    \includegraphics[width=\linewidth]{figs/fig_shrinkage.pdf}
    \caption{Multiplicative bias of the raw posterior means against the average posterior variance, in 15 equal-size bins of $\sigma_{\epsilon,i}$ on the Xu3 test set (6,667 galaxies per bin), for $g_1$ (blue circles) and $g_2$ (red squares). The dashed line is the shrinkage prediction of Eq.~(\ref{eqn:shrinkage}). The dotted vertical line marks the prior variance $\sigma_\pi^2$.}
    \label{fig:shrinkage}
\end{figure}

Fig.~\ref{fig:shrinkage} shows the result. The measured $m$ follows the prediction from the bins with intermediate posterior widths to $m \approx -1$ in the bins with the widest posteriors. Over the whole test set, the prediction gives $m_1 = -0.607$ and $m_2 = -0.617$, close to the measured $m_1 = -0.616 \pm 0.014$ and $m_2 = -0.592 \pm 0.014$. Shrinkage therefore explains almost all of the bias of the raw posterior means, which is dominated by galaxies with wide posteriors. This bias can be reduced by giving less weight to these galaxies (Sec.~\ref{ssec:weights-cuts}) or by calibrating the response of the estimator (Sec.~\ref{ssec:response}). A more natural approach, however, is to not reduce each posterior to a point estimate at all, and instead to combine the likelihoods of all galaxies that share the same shear, utilizing the full information in the posteriors (Sec.~\ref{ssec:stacking}).

\subsection{Applying Galaxy Weighting and Selection Cuts}
\label{ssec:weights-cuts}

As we saw in Sec.~\ref{ssec:shrinkage}, the bias of the raw posterior means is dominated by galaxies with wide posteriors. A straightforward way to reduce it is to give these galaxies less influence on the ensemble estimate, either by weighting galaxies according to their posterior widths or by removing the galaxies with the widest posteriors. Both are commonly done in WL. In this section we explore both, and compare cuts on posterior width with cuts on inclination, which were explored for DESI-KL by \citep{KLDESI}.

\subsubsection{Weighting}

We weight each galaxy by
\begin{equation}
\label{eqn:weights}
    w_i = \frac{1}{\sigma_{\epsilon,i}^2/2 + \sigma_\mathrm{floor}^2},
\end{equation}
where $\sigma_{\epsilon,i}^2/2$ is the average posterior variance of the two shear components of galaxy $i$, and refit $m$ and $c$ as in Sec.~\ref{ssec:raw-bias} with these weights. This has the same form as the inverse-variance weights used in WL, $w = 1/(\sigma_\mathrm{SN}^2 + \sigma_e^2)$, where $\sigma_\mathrm{SN}$ is the intrinsic shape noise and $\sigma_e$ is the measurement error of a galaxy \citep{HSCY1Shape}. In WL, the shape noise term is needed because the measurement error of a galaxy does not include the scatter of its intrinsic shape. Our posterior widths already include the full uncertainty of each galaxy, including its KL shape noise, so here the floor plays a different role of keeping the weights from depending too strongly on the posterior widths. We set it to the average posterior width per component over the test set, $\sigma_\mathrm{floor} = \langle\sigma_{\epsilon,i}\rangle/\sqrt{2} = 0.043$.

Without the floor, the weights $w_i \propto 1/\sigma_{\epsilon,i}^2$ differ by more than a factor of 10 between the 5\% of galaxies with the narrowest posteriors and the 5\% with the widest. The ensemble estimate is then dominated by a small number of galaxies: the effective number of galaxies, $N_\mathrm{eff} = (\sum_i w_i)^2/\sum_i w_i^2$, drops to 10.5\% of the test set. The result then relies on the posterior widths of these few galaxies being accurate individually, whereas our calibration test (Sec.~\ref{ssec:calibration}) only shows that they are accurate on average. With the floor, $N_\mathrm{eff}$ stays at 91\% of the test set.

The weights reduce the multiplicative bias from $m_1 = -0.616$ and $m_2 = -0.592$ to $m_1 = -0.512 \pm 0.016$ and $m_2 = -0.486 \pm 0.015$ with the floor, and to $m_1 = -0.210 \pm 0.019$ and $m_2 = -0.207 \pm 0.018$ without it. The additive bias stays consistent with zero in both cases. \XW{I've described the pros/cons of using shape noise floor vs not using it, as far as I understand it (also see fig \ref{fig:cuts}). I'm curious about other people's opinions on this.}

\subsubsection{Posterior Width and Inclination Cuts}

Instead of weighting, we can remove the galaxies with the widest posteriors, keeping only those with $\sigma_{\epsilon,i} < t$ for some threshold $t$. An alternative that has been considered for KL is to cut on inclination. Edge-on disks have larger shape noise in KL, and \citep{KLDESI} found that removing them improves the average shape noise of a sample, but argued that the loss in number density outweighs this gain. Our posteriors show the same trend: the median $\sigma_{\epsilon,i}$ is 0.035 for galaxies with $\cos i > 0.8$, and 0.079 for galaxies with $\cos i < 0.2$. Since the true inclination is not known for real galaxies, we cut on the mean of the TF-weighted posterior of $\cos i$ instead, keeping the galaxies that are closest to face-on. Both the posterior width and the posterior mean of $\cos i$ are available for every galaxy as soon as we run our NPE on its data.

\begin{figure}
    \includegraphics[width=\linewidth]{figs/fig_cuts.pdf}
    \caption{Top: average multiplicative bias $(m_1+m_2)/2$ of the posterior means on the Xu3 test set as a function of the fraction of galaxies kept, for cuts on the posterior width $\sigma_{\epsilon,i}$ (solid blue), on the posterior mean of $\cos i$ (dashed green) and on the true $\cos i$ (dotted green), all with equal weights. Shaded bands are the $\pm 1\sigma$ errors of the fit. Red markers show the weights of Eq.~(\ref{eqn:weights}) with and without the floor, plotted at their effective fraction $N_\mathrm{eff}/N$. Bottom: statistical uncertainty of the mean shear per component after dividing by the response, $\sigma_\mathrm{eff}/\sqrt{N_\mathrm{eff}}$ (Eq.~\ref{eqn:sigma-eff}), averaged over the two components, for the same cuts and weights.}
    \label{fig:cuts}
\end{figure}

Fig.~\ref{fig:cuts} compares the two cuts and the two weightings as a function of the fraction of galaxies kept. At the same fraction, the width cut reduces the bias much more than the inclination cut. Keeping the 32\% of galaxies with the narrowest posteriors gives an average $m = -0.11$, while keeping the 32\% closest to face-on gives $m = -0.27$. The inclination is only one of several properties that set the posterior width of a galaxy, along with its H$\alpha$ SNR, size and brightness, while the posterior width is directly the quantity that sets the amount of shrinkage (Eq.~\ref{eqn:shrinkage}). Cutting on the posterior mean of $\cos i$ gives almost the same result as cutting on the true $\cos i$, so the difference does not come from the uncertainty in the inclination. For the tightest width cuts, $m$ even becomes positive, reaching $m \approx +0.11$ when only 10\% of galaxies are kept, which shrinkage alone cannot explain. We further discuss this phenomenon later on in this section.

Removing or downweighting galaxies also changes the statistical uncertainty of the ensemble estimate, which we show in the bottom panel of Fig.~\ref{fig:cuts}. Since the different choices have very different $m$, we compare their uncertainties after correcting for the bias. We first correct each estimate using the $m_k$ and $c_k$ fitted on the test set, $\hat{g}^\mathrm{corr}_{k,i} = (\hat{g}_{k,i} - c_k)/(1+m_k)$, and then measure the spread of the corrected estimates around the true shear for the galaxies in the fit window, which gives an effective shape noise
\begin{equation}
\label{eqn:sigma-eff}
    \sigma_{\mathrm{eff},k} = \left[\frac{\sum_i w_i^2 \left(\hat{g}^\mathrm{corr}_{k,i} - g_{k,i}\right)^2}{\sum_i w_i^2}\right]^{1/2}.
\end{equation}
The uncertainty of the weighted mean of the corrected estimates is then $\sigma_{\mathrm{eff},k}/\sqrt{N_\mathrm{eff}}$. For a cut, $w_i$ is 1 for the kept galaxies and 0 for the removed ones. Since this uses the $m$ measured on simulations, which would not be known exactly for real data, the bottom panel of Fig.~\ref{fig:cuts} is only meant to show the trends between the different choices, and not necessarily the uncertainties we would obtain on real data.

For equal weights and no cut, $\sigma_\mathrm{eff}/\sqrt{N_\mathrm{eff}} = 1.79\times10^{-4}$ per component. Width cuts reduce it only slightly, to between $1.57\times10^{-4}$ and $1.74\times10^{-4}$, since removing galaxies reduces both the noise and the number of galaxies. Inclination cuts increase it, by up to 10\% for cuts on the posterior $\cos i$ and 30\% for cuts on the true $\cos i$. The weights with the floor give $1.57\times10^{-4}$, the same as the best width cuts, while keeping 91\% of the galaxies, and the weights without the floor give the smallest uncertainty of all, $1.08\times10^{-4}$. The weights without the floor therefore give both a smaller $|m|$ and a smaller statistical uncertainty than the weights with the floor, at the cost of relying on a small effective number of galaxies, and of a larger selection response (Sec.~\ref{sssec:selection}).

\subsubsection{Selection Bias}
\label{sssec:selection}

So far, we have treated the cuts and weights as if they only changed which galaxies contribute to the ensemble estimate. However, if any of these treatments depend on shear itself, that would induce a selection bias. To the first order, the response due to selection and the response due to shear are additive, such that
\begin{equation}
\label{eqn:selection-response}
    \langle\hat{g}\rangle_\mathrm{sel} \simeq (R_\gamma + R_\mathrm{sel})\, g, \qquad m = R_\gamma + R_\mathrm{sel} - 1,
\end{equation}
where $R_\gamma$ is the average response of the estimates of the selected galaxies, and $R_\mathrm{sel}$ is the extra response from our weighting or selection cuts. The posterior width responds to the shear, since, as we've discovered, galaxies whose measured shear is larger tend to have narrower posteriors (Sec.~\ref{ssec:shrinkage}). This causes galaxies with higher |g| to be preferentially selected by a width cut, pulling the mean shear amplitude up by selection alone, i.e. $R_\mathrm{sel} > 0$. Since shrinkage gives $R_\gamma < 1$, this extra response partly cancels the shrinkage bias, and for the tightest width cuts in Fig.~\ref{fig:cuts} it overcompensates, giving $m > 0$. A cut on the true $\cos i$, on the other hand, does not respond to the shear and has $R_\mathrm{sel} = 0$. Traditional WL measurements face the problem of selection bias, since their cuts on measured galaxy properties also respond to the shear, and methods like Metacalibration (Metacal) accounts for it by measuring $R_\mathrm{sel}$ directly and including it in the calibration \citep{MetacalPractical, Metadet}.

We measure $R_\mathrm{sel}$ in the same way as Metacal, on 1000 galaxies drawn from a catalog built like the Xu3 test set. Each galaxy is simulated at shears $(0,0)$, $(\pm\Delta, 0)$ and $(0, \pm\Delta)$ with $\Delta = 0.01$ and the same noise realization. Let $S_+$ be the selection obtained by applying the cut to the posterior widths of the versions sheared by $(+\Delta, 0)$, i.e. the set of galaxies whose sheared version passes the cut, and $S_-$ the same for the versions sheared by $(-\Delta, 0)$. Writing $\langle \hat{g}_1(0) \rangle_{S_\pm}$ for the mean of the estimates of the unsheared versions over the galaxies in $S_\pm$, the selection response is
\begin{equation}
\label{eqn:rsel}
    R_\mathrm{sel} = \frac{\langle \hat{g}_1(0) \rangle_{S_+} - \langle \hat{g}_1(0) \rangle_{S_-}}{2\Delta},
\end{equation}
and similarly for $g_2$. Since the same unsheared estimates are averaged in both terms, only the change in which galaxies are selected contributes. For weights, $S_\pm$ is replaced by the weights computed from the posterior widths of the sheared versions, and the means by weighted means. $R_\gamma$ is the average of the finite-difference responses $[\hat{g}_{1,i}(+\Delta) - \hat{g}_{1,i}(-\Delta)]/2\Delta$ of the selected galaxies.

\begin{table}
\centering
\caption{Selection response relative to the shear response, $R_\mathrm{sel}/R_\gamma$, averaged over the two shear components and measured on 1000 galaxies simulated at five shears with the same noise. Kept fractions refer to this sample. Errors are from bootstrap resampling of the galaxies.}
\label{tab:selection}
\begin{tabular}{lcc}
\toprule
Selection & Kept & $R_\mathrm{sel}/R_\gamma$ [\%] \\
\midrule
Width cut $\sigma_{\epsilon,i} < 0.05$ & 18\% & $+12 \pm 3$ \\
Width cut $\sigma_{\epsilon,i} < 0.075$ & 67\% & $+2.9 \pm 0.8$ \\
Weights with floor & -- & $+3.0 \pm 0.2$ \\
Weights without floor & -- & $+5.8 \pm 1.2$ \\
Posterior $\cos i > 0.8$ & 19\% & $+2.8 \pm 2.1$ \\
Posterior $\cos i > 0.6$ & 39\% & $-2.0 \pm 1.0$ \\
Posterior $\cos i > 0.44$ & 57\% & $-0.6 \pm 0.6$ \\
Posterior $\cos i > 0.31$ & 69\% & $-1.0 \pm 0.4$ \\
\bottomrule
\end{tabular}
\end{table}

Table~\ref{tab:selection} lists the results. The width cuts have a positive selection response that grows as the cut tightens, from $+2.9\%$ of $R_\gamma$ when 67\% of galaxies are kept to $+12\%$ when 18\% are kept. The selection response of the tighter cut, $R_\mathrm{sel} = 0.12$, is comparable to the difference in $m$ between width and inclination cuts at similar kept fractions in Fig.~\ref{fig:cuts}, so much of the advantage of the width cut in $m$ comes from selection rather than from reduced shrinkage. The weights have a smaller selection response, $+3.0\%$ with the floor and about twice that without it, while the cuts on the posterior $\cos i$ have selection responses within a few percent of zero. Any of these choices therefore leaves a bias that has to be calibrated, including its selection response, which we turn to next.

\subsection{Response Calibration}
\label{ssec:response}
\XW{The work in this section was done before I realized selection bias was a thing, so none of what I did here take that into account. Writing what I currently have in here for completeness reasons.}

Instead of reducing the influence of the galaxies with wide posteriors, we can try to correct for their bias. Eq.~(\ref{eqn:shrinkage}) predicts that the shear response $R = 1 + m$ of a galaxy depends on its posterior width, so if this relation $R(\sigma)$ is known, the estimates of an ensemble can be calibrated. Since $R(\sigma)$ cannot be measured on real data from our posteriors alone, we measure it on simulations, similar to how WL shape measurements are calibrated with image simulations \citep{DESY1Shape, HSCY1Sims}.

We draw 1000 galaxies from a catalog built in the same way as the Xu3 test set, and simulate each of them five times, at shears $(g_1, g_2) = (0, 0)$, $(\pm\Delta, 0)$ and $(0, \pm\Delta)$ with $\Delta = 0.01$, keeping all other parameters and the noise realization the same. For each galaxy, the shear response of each component is the finite difference
\begin{equation}
\label{eqn:response}
    R_{kk,i} = \frac{\hat{g}_{k,i}(+\Delta) - \hat{g}_{k,i}(-\Delta)}{2\Delta},
\end{equation}
where $\hat{g}_{k,i}(\pm\Delta)$ is the posterior mean of the version sheared by $\pm\Delta$ in component $k$. Since the two versions have the same noise, the noise largely cancels in this difference, and the response can be measured for each galaxy individually. We then fit $R_{kk,i} = a + b\,\sigma_{\mathrm{post},k,i}^2$ to both components together, using the posterior variances of the unsheared versions. The Gaussian model of Sec.~\ref{ssec:shrinkage} predicts $a = 1$ and $b = -1/\sigma_\pi^2$.

\begin{figure*}
    \includegraphics[width=\linewidth]{figs/fig_response.pdf}
    \caption{Shear response of the posterior means, measured on 1000 galaxies drawn like the Xu3 test set, each simulated at five shears with the same noise. (a) Response $R_{kk}$ of each galaxy and component (grey points) against its posterior variance, with means in deciles of the posterior variance (circles), the linear fit $a + b\,\sigma_\mathrm{post}^2$ (solid) and the Gaussian prediction of Eq.~(\ref{eqn:shrinkage}) (dashed). (b) Mean response in bins of the true $\cos i$, with the Gaussian prediction from the mean posterior variance in each bin. Error bars are from bootstrap resampling of the galaxies.}
    \label{fig:response}
\end{figure*}

Fig.~\ref{fig:response}a shows the result. The response falls with the posterior variance as predicted, with a fitted slope similar to the prediction. The fitted intercept is above 1, as the galaxies with the narrowest posteriors have a response slightly above 1. Fig.~\ref{fig:response}b shows the same response as a function of inclination: nearly edge-on galaxies respond very little to the shear, and the response increases toward face-on galaxies, which have the narrowest posteriors.

To apply the fitted $R(\sigma)$ to the Xu3 test set, the simplest approach would be to divide the estimate of each galaxy by its predicted response, $\hat{g}_{k,i}/R(\sigma_{k,i})$. However, this is unstable for the many galaxies with $R$ close to zero, whose noise is amplified by $1/R$, and leaves a large bias. We instead calibrate the ensemble by fitting $\hat{g}_{k,i} = c_k + (1+m_k)\,R(\sigma_{k,i})\,g_{k,i}$ over the fit window, which is equivalent to averaging the corrected estimates $\hat{g}_{k,i}/R(\sigma_{k,i})$ with weights $R(\sigma_{k,i})^2$, so that galaxies with a small response get little weight. Galaxies with a predicted response below zero get zero weight. This gives $m_1 = -0.049 \pm 0.029$ and $m_2 = -0.000 \pm 0.030$, with the additive bias consistent with zero, where the errors come from bootstrap resampling of both the 1000 calibration galaxies and the test set. The effective number of galaxies is 42\% of the test set, and the effective shape noise (Eq.~\ref{eqn:sigma-eff}) gives $\sigma_\mathrm{eff}/\sqrt{N_\mathrm{eff}} = 1.34\times10^{-4}$ per component, lower than for any of the weights or cuts of Sec.~\ref{ssec:weights-cuts} except the weights without the floor.
